\section{Trabalhos Relacionados}
\label{sec:related}

Para a construção do referencial teórico deste trabalho, realizou-se uma análise sistemática de publicações recentes que investigam a segurança de prompts em Modelos de Linguagem de Grande Escala (LLMs). Os estudos examinados abrangem propostas de frameworks de defesa, sistematizações de conhecimento, mecanismos baseados em consultas estruturadas, análises empíricas de vulnerabilidades e avaliações voltadas à privacidade e percepção dos usuários. Essa revisão permitiu delinear o estado da arte, evidenciando o progresso alcançado na compreensão e mitigação de ataques de injeção de prompt e ameaças correlatas.

Entretanto, a análise comparativa revelou que cada um desses trabalhos apresenta limitações relevantes. Como pode ser observado na Tabela 1:

\small
\setlength{\extrarowheight}{2pt}

\begin{longtable}{
      >{\raggedright\arraybackslash}p{0.22\textwidth}
      >{\raggedright\arraybackslash}p{0.18\textwidth}
      >{\raggedright\arraybackslash}p{0.18\textwidth}
      >{\raggedright\arraybackslash}p{0.18\textwidth}
      >{\raggedright\arraybackslash}p{0.20\textwidth}}
\caption{Comparison of Related Works on Prompt Injection Defense in LLMs}
\label{tab:related_works}
\\
\toprule
\textbf{Referencia} & \textbf{Foco} & \textbf{Metodologia} & \textbf{Métricas} & \textbf{Lacuna} \\
\midrule
\endfirsthead

      Lin, Huawei \emph{et al.} \cite{Lin2025UniGuardian} 
      & Proposta de defesa unificada (UniGuardian) contra ataques de LLM (injeção, backdoor, adversariais).
      & Propõe um framework de detecção unificado para analisar prompts e saídas.
      & (1) auROC (Area Under the Receiver Operator Characteristic Curve) e (2) auPRC (Area Under the Precision-Recall Curve).
      & Falta de profundidade no que tange a ataques de backdoor mais complexos ou ofuscados, gerando falsos positivos. \\
    \midrule
      Hong, Hanbin \emph{et al.} \cite{Hong2025SystematizationOfKnowledge} 
      & Sistematização do conhecimento em segurança de prompts em LLMs, propondo uma taxonomia e um conjunto de métricas próprias para padronizar avaliações de ataques e defesas.
      & Revisão sistemática e proposta de uma taxonomia multi-nível. Desenvolvimento de toolkit aberto e dataset (JailbreakDB) para avaliação padronizada.
      & Propõe métricas padronizadas próprias integradas a um toolkit de avaliação com taxonomia hierárquica e perfis de ameaça formais.
      & Limitação na aplicação prática das métricas e taxonomias propostas; necessidade de validação contínua e ampliação para modelos multimodais e cenários reais. \\
    \midrule
      Chen, Sizhe \emph{et al.} \cite{Chen2025StruQ}
      & Proposta de defesa chamada \textit{StruQ}, que utiliza consultas estruturadas (\textit{structured queries}) para separar dados do prompt, reduzindo o risco de injeções maliciosas em LLMs.
      & Implementação experimental com análise de desempenho e eficácia em diferentes cenários de injeção. Avalia a separação entre dados e instruções para mitigar ataques.
      & Taxa de sucesso de ataque, latência e sobrecarga do sistema (\textit{defense overhead}).
      & Boa eficácia em prompts simples, mas limitação em ataques indiretos e cenários complexos de múltiplas etapas, exigindo integração com outras técnicas de defesa. \\
      
    \midrule
      Benjamin, Victoria \emph{et al.} \cite{benjamin2024}
      & Análise sistemática da vulnerabilidade de 36 LLMs a ataques de injeção de prompt focados em gerar código de keylogger.
      & 4 prompts de injeção direta contra 36 LLMs. Análise estatística (Correlação, Random Forest, SHAP, PCA).
      & (1) Taxa de sucesso (2) Importância de features. (3) Correlação entre prompts.
      & Necessidade de testes multilíngues, múltiplas etapas de complexidade. \\
    \midrule
      Sebastian, Glorin. \cite{Sebastian2023SecurityUserInformation:}
      & Investigar a proteção de dados e privacidade em chatbots (foco no ChatGPT). Avaliar Tecnologias de Aprimoramento de Privacidade (PETs).
      & Revisão da literatura, análise de técnicas (ex: privacidade diferencial) e uma pesquisa (survey) com 177 usuários.
      & Métricas de percepção da pesquisa: (1) Nível de preocupação, (2) Disposição para sacrificar desempenho, (3) Consciência sobre vazamentos.
      & Análise de riscos à privacidade, preocupação dos usuários através de pesquisas. \\
    \midrule
    \textbf{This Work}
      & \textbf{Análise comparativa de frameworks de defesa contra injeção de prompt em LLMs.}
      & \textbf{Prova de conceito e simulações baseadas em diferentes cenários de ataque.}
      & \textbf{Taxa de sucesso de ataque, resiliência e custo computacional.}
      & \textbf{Busca preencher a lacuna de falta de metodologia padronizada para comparação empírica entre defesas.} \\
    \bottomrule
\end{longtable}

Nesse sentido, as lacunas observadas justificam a necessidade de investigações adicionais que consolidem métricas, experimentos e metodologias, oferecendo uma base mais consistente para o desenvolvimento e comparação de soluções de segurança em LLMs.